% !TeX root = ../Tempel-Daniel.tex
\chapter{Introduction}
\label{chap:introduction}

Large language models (LLMs) are increasingly used as components of software development tools. Modern AI-assisted development systems can analyse source code, explain errors, generate modifications and interact with external tools such as version-control systems, test runners and documentation sources. However, building such a system requires more than selecting a capable language model. Model deployment, context management, tool integration, retrieval, security and evaluation all influence whether the resulting system is practical and reliable.

This report examines these topics through a common example: an \textbf{AI Software Engineering Assistant} whose task is to analyse a GitHub issue, inspect the relevant source code, propose a modification and verify the solution by running tests. The same use case is developed throughout the report so that the individual AI concepts can be considered as parts of one system rather than as isolated technologies.

The report begins with the foundations of running large language models locally and the hardware constraints this creates. It then considers how models can be compared and selected for a software-engineering workload. Later sections extend the assistant with structured outputs and tool use, context management, retrieval, the Model Context Protocol and autonomous agent behaviour. Finally, production architecture, security and evaluation are considered to examine how such a system could be operated reliably.

The aim is not to implement a complete production assistant, but to understand the main architectural concepts and trade-offs involved in constructing modern AI systems for software engineering.
